
\documentclass{article}
\usepackage{colortbl}
\usepackage{makecell}
\usepackage{multirow}
\usepackage{supertabular}

\begin{document}

\newcounter{utterance}

\centering \large Interaction Transcript for game `wordle', experiment `high\_frequency\_words\_no\_clue\_no\_critic', episode 10 with reflexion{-}Qwen3.5{-}9B.
\vspace{24pt}

{ \footnotesize  \setcounter{utterance}{1}
\setlength{\tabcolsep}{0pt}
\begin{supertabular}{c@{$\;$}|p{.15\linewidth}@{}p{.15\linewidth}p{.15\linewidth}p{.15\linewidth}p{.15\linewidth}p{.15\linewidth}}
    \# & \multicolumn{2}{c}{Player} && \multicolumn{2}{c}{Game Master} \\
    \hline

    \theutterance \stepcounter{utterance}  
    & & & \multicolumn{4}{p{0.6\linewidth}}{
        \cellcolor[rgb]{0.9,0.9,0.9}{
            \makecell[{{p{\linewidth}}}]{
                \texttt{\tiny{[P1$\langle$GM]}}
                \texttt{You are a language wizard who likes to guess words by using the given rules.} \\
\\ 
\texttt{Welcome to Wordle! You have six attempts to guess the target word, a valid English word of five lowercase letters (a{-}z). Please use the tags "explanation:" and "guess:" to provide a concise explanation for each guess.} \\
\\ 
\texttt{For instance, if your guess is "apple", your response should be} \\
\texttt{explanation: this is a common five{-}letter English word, and I am starting my guess with this word.} \\
\texttt{guess: apple} \\
\\ 
\texttt{After each guess, your answer will be validated, and you will receive feedback indicating which letters are correct (green), which letters are correct but in the wrong position (yellow), and which letters are incorrect (red). This feedback can be useful in determining which letters to include or exclude in your next guess.} \\
\\ 
\texttt{For example, the feedback for "apple" might be:} \\
\texttt{guess\_feedback: a<yellow> p<yellow> p<green> l<yellow> e<red>} \\
\\ 
\texttt{The explanation should contain details about how the guess\_feedback is used to arrive at a new guess.} \\
\\ 
\texttt{Do not generate any other text.} \\
\\ 
\texttt{Let's begin with your first guess.} \\
            }
        }
    }
     \\ \\

    \theutterance \stepcounter{utterance}  
    & \multicolumn{4}{p{0.6\linewidth}}{
        \cellcolor[rgb]{0.561,0.737,0.561}{
            \makecell[{{p{\linewidth}}}]{
                \texttt{\tiny{[P1$\rangle$GM]}}
                \texttt{explanation: I am starting with "crane" because it contains common vowels (a, e) and frequent consonants (r, n) that appear in many five{-}letter words, allowing me to quickly eliminate or confirm several letters.} \\
\texttt{guess: crane} \\
            }
        }
    }
    & & \\ \\

    \theutterance \stepcounter{utterance}  
    & & & \multicolumn{2}{p{0.3\linewidth}}{
        \cellcolor[rgb]{0.9,0.9,0.9}{
            \makecell[{{p{\linewidth}}}]{
                \texttt{\tiny{[GM$|$GM]}}
                \texttt{attempts = 1} \\
\texttt{target\_word = cover} \\
\texttt{guess = crane} \\
\texttt{guess\_feedback = c<green> r<yellow> a<red> n<red> e<yellow>} \\
            }
        }
    }
    & & \\ \\

    \theutterance \stepcounter{utterance}  
    & & & \multicolumn{4}{p{0.6\linewidth}}{
        \cellcolor[rgb]{0.9,0.9,0.9}{
            \makecell[{{p{\linewidth}}}]{
                \texttt{\tiny{[P1$\langle$GM]}}
                \texttt{guess\_feedback: c<green> r<yellow> a<red> n<red> e<yellow>} \\
\\ 
\texttt{Provide your response only in this format.} \\
\texttt{explanation: details} \\
\texttt{guess: word} \\
            }
        }
    }
     \\ \\

    \theutterance \stepcounter{utterance}  
    & \multicolumn{4}{p{0.6\linewidth}}{
        \cellcolor[rgb]{0.561,0.737,0.561}{
            \makecell[{{p{\linewidth}}}]{
                \texttt{\tiny{[P1$\rangle$GM]}}
                \texttt{explanation: The feedback 'c<green> r<yellow> a<red> n<red> e<yellow>' confirms 'c' is the first letter, 'r' and 'e' are present but misplaced, while 'a' and 'n' are absent. Since 'r' is second in 'crane' but yellow, it must be elsewhere; 'e' is last in 'crane' but yellow, so it must be elsewhere. I will keep 'c' at the start, try placing 'e' in the middle (position 3) and 'r' in position 4, and replace 'a' and 'n' with other common letters like 'o' and 't' to form a new candidate.} \\
\texttt{guess: coter} \\
            }
        }
    }
    & & \\ \\

    \theutterance \stepcounter{utterance}  
    & & & \multicolumn{2}{p{0.3\linewidth}}{
        \cellcolor[rgb]{0.9,0.9,0.9}{
            \makecell[{{p{\linewidth}}}]{
                \texttt{\tiny{[GM$|$GM]}}
                \texttt{Error: The guessed word is not a valid word for this game.} \\
            }
        }
    }
    & & \\ \\

    \theutterance \stepcounter{utterance}  
    & & & \multicolumn{4}{p{0.6\linewidth}}{
        \cellcolor[rgb]{0.9,0.9,0.9}{
            \makecell[{{p{\linewidth}}}]{
                \texttt{\tiny{[P1$\langle$GM]}}
                \texttt{Your guess is not a valid word for this game. Please try again.} \\
\\ 
\texttt{Provide your response only in this format.} \\
\texttt{explanation: details} \\
\texttt{guess: word} \\
            }
        }
    }
     \\ \\

    \theutterance \stepcounter{utterance}  
    & \multicolumn{4}{p{0.6\linewidth}}{
        \cellcolor[rgb]{0.561,0.737,0.561}{
            \makecell[{{p{\linewidth}}}]{
                \texttt{\tiny{[P1$\rangle$GM]}}
                \texttt{explanation: The previous guess "coter" was invalid. Based on the feedback from "crane" (c<green> r<yellow> a<red> n<red> e<yellow>), I know the word starts with 'c', contains 'r' and 'e' but not 'a' or 'n'. 'r' is not in the second position, and 'e' is not in the fifth. I will try "creed" to test if 'r' is in the third position and 'e' is in the fourth, while keeping 'c' at the start.} \\
\texttt{guess: creed} \\
            }
        }
    }
    & & \\ \\

    \theutterance \stepcounter{utterance}  
    & & & \multicolumn{2}{p{0.3\linewidth}}{
        \cellcolor[rgb]{0.9,0.9,0.9}{
            \makecell[{{p{\linewidth}}}]{
                \texttt{\tiny{[GM$|$GM]}}
                \texttt{attempts = 2} \\
\texttt{target\_word = cover} \\
\texttt{guess = creed} \\
\texttt{guess\_feedback = c<green> r<yellow> e<red> e<green> d<red>} \\
            }
        }
    }
    & & \\ \\

    \theutterance \stepcounter{utterance}  
    & & & \multicolumn{4}{p{0.6\linewidth}}{
        \cellcolor[rgb]{0.9,0.9,0.9}{
            \makecell[{{p{\linewidth}}}]{
                \texttt{\tiny{[P1$\langle$GM]}}
                \texttt{guess\_feedback: c<green> r<yellow> e<red> e<green> d<red>} \\
\\ 
\texttt{Provide your response only in this format.} \\
\texttt{explanation: details} \\
\texttt{guess: word} \\
            }
        }
    }
     \\ \\

    \theutterance \stepcounter{utterance}  
    & \multicolumn{4}{p{0.6\linewidth}}{
        \cellcolor[rgb]{0.561,0.737,0.561}{
            \makecell[{{p{\linewidth}}}]{
                \texttt{\tiny{[P1$\rangle$GM]}}
                \texttt{explanation: The feedback 'c<green> r<yellow> e<red> e<green> d<red>' confirms the word starts with 'c', has 'e' as the fourth letter, and contains 'r' and another 'e' (which is already placed at position 4, so the yellow 'r' must be elsewhere, and the red 'e' at position 3 means the second 'e' in my guess was a false positive or the word only has one 'e' at position 4? Wait, the feedback shows two 'e's in the guess: position 3 is red, position 4 is green. This means the word has an 'e' at position 4, but NOT at position 3. The 'r' is yellow, meaning it exists but is not at position 2. Since 'c' is fixed at 1 and 'e' is fixed at 4, and 'r' is not at 2, 'r' must be at 3 or 5. However, position 3 was 'e' in my guess and got red, so 'r' could be at 3. Let's try placing 'r' at position 3. The letters 'd' and the second 'e' (if any) are incorrect or misplaced. I will guess "creed" again? No, that was the guess. I need a new word. If 'r' is at 3, the pattern is c{-}r{-}e{-}e{-}?. But wait, the feedback said 'e' at pos 3 is red. So the word is c{-}?{-}e{-}?{-}?. 'r' is yellow, so it's in the word but not at pos 2. It could be at 3, 5. If I put 'r' at 3, then pos 3 is 'r', not 'e'. But the feedback for 'creed' said pos 3 'e' is red. That means the word does NOT have 'e' at pos 3. It DOES have 'e' at pos 4. So the word is c{-}?{-}?{-}e{-}?. 'r' is yellow, so it's in the word. If 'r' is at 3, then the word is c{-}?{-}r{-}e{-}?. If 'r' is at 5, then c{-}?{-}?{-}e{-}r. Let's try "crane" again? No, 'a' and 'n' were red. Let's try "creek". c{-}r{-}e{-}e{-}k. Pos 2 'r' is yellow (ok), pos 3 'e' is red (ok, word doesn't have e at 3), pos 4 'e' is green (ok), pos 5 'k' is red. This fits perfectly if the word is "creek" but wait, "creek" has} \\
            }
        }
    }
    & & \\ \\

    \theutterance \stepcounter{utterance}  
    & & & \multicolumn{2}{p{0.3\linewidth}}{
        \cellcolor[rgb]{0.9,0.9,0.9}{
            \makecell[{{p{\linewidth}}}]{
                \texttt{\tiny{[GM$|$GM]}}
                \texttt{Error: The response should contain the 'guess:' keyword exactly once.} \\
            }
        }
    }
    & & \\ \\

    \theutterance \stepcounter{utterance}  
    & & & \multicolumn{2}{p{0.3\linewidth}}{
        \cellcolor[rgb]{0.9,0.9,0.9}{
            \makecell[{{p{\linewidth}}}]{
                \texttt{\tiny{[GM$|$GM]}}
                \texttt{game\_result = ABORT} \\
            }
        }
    }
    & & \\ \\

    \theutterance \stepcounter{utterance}  
    & & \multicolumn{4}{p{0.6\linewidth}}{
        \cellcolor[rgb]{0.95,0.95,0.95}{
            \makecell[{{p{\linewidth}}}]{
                \texttt{\tiny{[P1$\leftrightarrow$P]}}
                \texttt{Post{-}episode Reflexion (not sent to GM):} \\
\texttt{New plan: The previous attempt was aborted because my explanation became excessively long and rambling, violating the instruction to provide a "concise" explanation. Additionally, I failed to finalize a valid guess within the explanation before the text cut off. Next time, I will strictly limit the explanation to one or two sentences that directly map the feedback to the specific letter positions of the new word, ensuring the resulting word is a valid five{-}letter English word that adheres to all constraints (e.g., 'c' at start, 'e' at position 4, 'r' not at position 2) before submitting the "guess:" tag.} \\
            }
        }
    }
    & \\ \\

\end{supertabular}
}

\end{document}
